\section{Intervention and Participant Discontinuation/Withdrawal}

\draftinstr{Follow the Phase 2 publication section
\texttt{final-protocol/publication/sections/sec-07-discontinuation.tex} and the
Phase 1 model file
\texttt{trial-protocol/draft-protocol/sections/sec-07-discontinuation.tex}. No
figure or table is rendered in this section. Distinguish discontinuation of an
intervention from withdrawal from the study: a participant who stops one or both
components for a protocol-defined reason remains in the randomized ITT population
for survival follow-up. Apply rules within the randomized arm. Source from
\texttt{trial-protocol/nih-protocol/05\_intervention\_and\_participant\_discontinuation.md}.}

\subsection{Discontinuation of Study Intervention}
\draftinstr{Device stopping rules (Arm A), implementing 21 CFR \S312.23(g)(i)(v)
and the divergence-reporting threshold of \S312.32(g) \cite{kawchak2026adapt312}:
conversion to open or manual technique; repeated force-limit excursions against the
$\leq 3$ N / $\leq 18$ N caps; no-fly-zone violations (SMV/PV 1.0 mm; HA/CA/SMA
1.5 mm); sim-to-real divergence beyond twice tolerance ($> 3$ mm or $> 0.8$ N);
and any hard-stop e-stop. Drug stopping rules: daraxonrasib discontinued for
unresolved toxicity, confirmed ISGPS Grade C fistula \cite{Bassi2017}, pregnancy,
contraindicating illness, or request; modified FOLFIRINOX managed per standard
toxicity rules \cite{Conroy2018}. The federated hash-chained audit trail is
preserved across each event (\S10) and reviewed by the Physical AI Safety Review
Committee and the DSMB before further device use.}

\subsection{Participant Discontinuation/Withdrawal from the Study}
\draftinstr{Voluntary withdrawal at any time without penalty; separate withdrawal
of consent for procedures while permitting survival follow-up; investigator-
initiated withdrawal for safety or non-compliance. Data and specimens collected up
to withdrawal are retained. Because this is a randomized ITT trial, participants are
analyzed as randomized and are not replaced (\S9).}

\subsection{Lost to Follow-Up}
\draftinstr{Define lost-to-follow-up only after at least three documented telephone
attempts across separate days plus a written attempt over the 24-month OS window,
with vital-status ascertainment from the record, referring provider, and
permissible public sources; record the last date known alive. Missing-data and
censoring handling is specified in \S9.}
